\documentclass[a4paper,12pt]{article}
\usepackage[UTF8]{ctex}
\usepackage{amsmath,amssymb}
\usepackage{geometry}
\geometry{left=2.5cm,right=2.5cm,top=2.5cm,bottom=2.5cm}
\usepackage{longtable}
\usepackage{array}
\usepackage{hyperref}
\hypersetup{hidelinks}

\setlength{\parindent}{0pt}

\begin{document}
% Force cache refresh

\begin{center}
    {\Large \textbf{\underline{信息工程学院实践教学环节}}}\\[1em]
    {\Large \textbf{\underline{《面向对象的图表生成系统》实施计划书}}}
\end{center}

\vspace{1em}

\textbf{姓名：}吉天睿 \hfill \textbf{学号：}2025013446 \hfill \textbf{班级：}计算机类2505 \\[0.5em]
\textbf{邮箱：}tianruiji@nwafu.edu.cn

\vspace{1em}

\renewcommand{\arraystretch}{1.5}
\begin{longtable}{|>{\centering\arraybackslash}p{2.5cm}|p{4.7cm}|>{\centering\arraybackslash}p{2.5cm}|p{\dimexpr\textwidth-9.7cm-8\tabcolsep-5\arrayrulewidth\relax}|}
\hline
\textbf{项目名称} & \multicolumn{3}{p{\dimexpr\textwidth-2.5cm-6\tabcolsep-3\arrayrulewidth\relax}|}{面向对象的图表生成系统} \\
\hline
\textbf{指导教师} & 赵研 & \textbf{电子邮件} & yanzhao@nwafu.edu.cn \\
\hline
\textbf{任务描述} & \multicolumn{3}{p{\dimexpr\textwidth-2.5cm-6\tabcolsep-3\arrayrulewidth\relax}|}{%
\textbf{一、总体任务：}\newline
1. 数据加载：从CSV文件读取数据转化为图表\newline
2. 图表绘制：实现柱状图、饼图、折线图三种基本图表类型，并从折线图派生出面积图。通过抽象基类统一接口。\newline
3. 数据排序：支持按数值升序/降序、按名称升序/降序排序，并且支持重置为原始数据。\newline
4. 数据统计：计算并显示数据条数、最大值、最小值、平均值、中位数。\newline
5. 图片导出：支持单个图表导出和多类型图表批量导出为PNG文件。\newline
6. 图形界面：基于EasyX实现主页面和图表页面，包含按钮、输入框、弹窗等交互控件。\newline
7. 主题切换：提供六套颜色主题，图表颜色与界面同步切换。
\vspace{0.5em}\newline
\textbf{二、个人任务：}\newline
1. 主程序流程设计\newline
2. 页面类体系设计\newline
3. 基础界面控件实现\newline
4. 导出路径对话框\newline
5. 主题切换与界面美化\newline
6. 项目运行配置
} \\
\hline
\textbf{技术路线} & \multicolumn{3}{p{\dimexpr\textwidth-2.5cm-6\tabcolsep-3\arrayrulewidth\relax}|}{%
1. 按需分析与模块划分：分析项目功能需求，确定界面层与图表层的分工。\newline
2. 类图设计：设计Widget、Button、Card、TextInput、Page等类的继承关系和接口。\newline
3. 控件实现：实现Widget基类和派生类Button、Card、TextInput、PopupCard。\newline
4. 页面实现：实现Page基类和派生类MainPage、ChartPage，集成控件并绑定回调函数。\newline
5. 主程序集成：编写main.cpp，初始化窗口、创建页面、实现事件循环。\newline
6. 主题系统实现：定义ColorTheme结构体，实现主题切换逻辑。\newline
调试与测试：测试各控件交互、页面切换、主题切换、导出功能。\newline
7. 使用VSCode作为代码编辑器，使用EasyX作为基础图形库，CPP作为编程语言进行实现。
\vspace{0.5em}\newline
下面给出整体流程图：
} \\
\hline
\textbf{进度安排} & \multicolumn{3}{p{\dimexpr\textwidth-2.5cm-6\tabcolsep-3\arrayrulewidth\relax}|}{%
6月14日：选题\newline
6月15日：项目总体设计与模块划分\newline
6月16日：设计类图\newline
6月17日---20日：编写代码，中期检查\newline
6月21日---23日：编写代码，程序调试\newline
6月24日---25日：撰写实习报告，制作幻灯片准备答辩\newline
6月26日：实习答辩
} \\
\hline
\textbf{考核要求} & \multicolumn{3}{p{\dimexpr\textwidth-2.5cm-6\tabcolsep-3\arrayrulewidth\relax}|}{%
总成绩 ＝ 考勤成绩 × 20\% + 答辩成绩 × 30\% + 报告成绩 × 50\%
} \\
\hline
\end{longtable}

\vspace{2em}
\hfill \textbf{日期：} 2026年6月17日

\end{document}
